\section{合卷：改革扩大了什么，又不能替代什么}

新疆设省、台湾建省、海军与铁路、学堂与新军、修律与自治，使晚清国家拥有比过去更多的制度语言和技术渠道。甲午战败、赔款、租借和日俄战争又使这些项目必须在国际财政和军事竞争中推进。改革不是虚词；它改变了资格、组织、信息和地方争论的空间。

然而，设局不等于有经费，章程不等于有人员，学校与警察的出现也不等于居民拥有平等的进入机会。边地、灾区、港口、城市和乡村的节奏不同；旧有宗族、州县、商会和军队并未因新名词而退出。新政的矛盾不在于“改革或保守”的二分，而在于国家需要地方合作和财政资源，却又试图以更集中的方式取得它们。

\subsection{制度得失}

\textbf{所得：} 教育、军政、法律与公共讨论的制度框架扩展，新的职业与政治行动者得以出现。

\textbf{代价：} 财政、土地、身份和代表资格的门槛没有同步消失，外债和铁路尤其使中央与地方的利益冲突公开化。

\textbf{留下的压力：} 立宪机构与新式传播将把这些冲突带入1911年的革命政治。
